\section{Conclusion et perspectives}

\subsection{Ce qu'il faut retenir}

\begin{retenir}
\begin{enumerate}[nosep,leftmargin=*]
  \item \textbf{L'AD donne la tangente cohérente exacte.} Elle coïncide à $3\times10^{-16}$ près avec la formule de Simo--Hughes, et les différences finies plafonnent à $10^{-11}$. Écrire $\psi$ ou \code{integrate} suffit : la dérivation à la main devient un \emph{test}, et n'est plus une étape du développement.
  \item \textbf{Trois pièges, trois tests.} Un \code{if} est remplacé par \code{where}. Un \code{where} seul laisse passer des NaN en mode inverse, d'où le double \code{where}. La dérivation par rapport à une matrice $3\times3$ casse les symétries mineures, d'où le passage en Mandel.
  \item \textbf{Bien choisir le cas test.} Écrouissage linéaire et chargement proportionnel donnent une réponse affine par morceaux et un Newton qui converge en une itération. Pour \emph{voir} la convergence quadratique, il faut un écrouissage non linéaire (Voce) ou un chargement non proportionnel. On mesure alors un ordre de 2,01 contre 1,00 avec la tangente élastique, dont le taux de 0,130 se prédit par le spectre de $\mathbf{I}-\mathbf{K}_e^{-1}\mathbf{K}_t$.
  \item \textbf{Le noyau est limité par la mémoire.} JAX CPU est 3,5 à 4,5 fois plus rapide que NumPy grâce à la fusion. XLA laisse cependant un trafic 3,5 fois supérieur au minimum, qu'un noyau écrit à la main pourrait éliminer.
  \item \textbf{Le GPU n'a d'intérêt que si les données y résident.} Mesuré sur une RTX~2050 : avec données résidentes, il est 2 fois plus rapide que le CPU en FP64 et 4 fois en FP32. Avec transferts, il est 2 à 3 fois plus lent, car le bus PCIe prend près de 90\,\% du temps. Le FP64 de la carte vaut 1/52 de son FP32.
  \item \textbf{Un réseau peut être physique.} Un ICNN de 230 paramètres, polyconvexe par construction, reproduit les trois essais de Treloar à 3--6\,\% près et prédit le cisaillement pur non vu à 12\,\% près, contre 21\,\% pour le meilleur polynôme.
\end{enumerate}
\end{retenir}

\subsection{Perspectives}

\begin{itemize}[nosep]
  \item \textbf{Garder l'état sur le GPU} entre les itérations de Newton, et ne faire circuler que $\eps$ et $\sig$ : le plafond PCIe passerait de $10^7$ à $3\times10^7$ points/s en FP64, au-dessus du CPU (section~\ref{sec:gpu}). Utiliser aussi la mémoire hôte verrouillée pour accélérer le retour vers l'hôte.
  \item \textbf{Noyau fusionné} (Pallas ou CUDA) pour approcher le trafic minimal de 304~octets par point, et mesurer le gain par rapport à XLA.
  \item \textbf{Couplage avec un code EF}, par exemple avec la tangente exportée vers code\_aster ou FEniCSx, ou un assemblage complet en JAX pour garder l'état sur GPU.
  \item \textbf{Lois plus riches} : écrouissage cinématique non linéaire (Armstrong--Frederick, Chaboche), viscoplasticité, plasticité en grandes déformations. La dérivation implicite de la section~\ref{sec:implicit} s'étend telle quelle à un système non linéaire vectoriel.
  \item \textbf{Identification de paramètres} par gradient en mode inverse à travers \code{integrate}, sur une courbe d'essai complète. C'est le cas d'usage où le double \code{where} devient indispensable.
  \item \textbf{Précision mixte} : calculer la loi en FP32 et le résidu d'équilibre en FP64, en mesurant l'effet sur la convergence du Newton global.
\end{itemize}
